% ============================================================
%  5. ARQUITECTURA FÍSICA DE LA SOLUCIÓN  (Subdocumento 4.2 — Formulario T-11)
%  Informe 1 — Arquitectura física
%  a) Arquitectura Física · b) Tecnologías de software · c) Implementos a proveer
%  d) Data Center primario · e) Data Center secundario
% ============================================================
\chapter{Arquitectura Física de la Solución}
\label{sec:arqfisica}

Este capítulo describe la arquitectura física de la solución, incluido el modelo de
emplazamiento híbrido obligatorio, las tecnologías de software, el equipamiento a
proveer, y las especificaciones de los data centers primario y secundario. El detalle se
complementa en el Formulario T-11.

\section{Modelo de emplazamiento híbrido}

La solución es obligatoriamente híbrida (nube pública regional + componentes
\emph{on-premise} críticos), conforme al Artículo 16° de las Bases Administrativas y a
RT-03.01. El modelo de emplazamiento componente por componente es:

\begin{table}[htbp]
\centering
\small
\caption{Modelo de emplazamiento híbrido}
\label{tab:emplazamiento}
\begin{tabularx}{\textwidth}{@{}L{5.4cm}L{4.6cm}X@{}}
\toprule
\textbf{Componente} & \textbf{Emplazamiento} & \textbf{Justificación} \\
\midrule
Recepción, \emph{picking} y carga
(CD Talca y Concepción) &
On-premise, con réplica a nube &
Debe continuar operando durante un corte del enlace (Restricción N° 2 del Caso);
ambiente de bodega con montacargas y cámara de congelado a $-$22\,°C sin cobertura de
señal. \\
\addlinespace
Plataformas de \emph{cross-docking}
(Curicó, Chillán, Los Ángeles) &
Gabinete de borde; conectividad exclusiva por red móvil &
Sin almacenamiento propio, ventana de operación de tres horas, sin fibra óptica
disponible. \\
\addlinespace
Preventa y reparto (dispositivos
móviles de terreno) &
Aplicación \emph{offline-first}, sincronización diferida a nube &
Deben operar un turno completo sin señal, en rutas rurales de hasta 240 km
(Restricción N° 3 del Caso). \\
\addlinespace
Trazabilidad de lote, cadena de frío,
analítica, portal, componentes de IA &
Nube pública (región Chile o Sudamérica) &
Volumen variable y sin acoplamiento a equipamiento físico; elasticidad estacional
durante el \emph{peak} de septiembre. \\
\addlinespace
Telemetría de los 42 camiones propios &
Se integra desde la plataforma del proveedor externo que ya la provee &
Fuente existente a preservar; los 54 camiones de transportistas no están cubiertos y se
incorporan mediante mecanismo a definir con cada empresa (Decisión N° 5). \\
\bottomrule
\end{tabularx}
\end{table}

\section{Tecnologías de software}

Las tecnologías de software se eligen bajo el criterio de neutralidad tecnológica y
vigencia de soporte por los 56 meses contractuales (Restricción R-29):

\begin{table}[htbp]
\centering
\small
\caption{Tecnologías de software propuestas}
\label{tab:software}
\begin{tabularx}{\textwidth}{@{}L{5.2cm}X@{}}
\toprule
\textbf{Componente} & \textbf{Tecnología propuesta} \\
\midrule
Proveedor de nube pública &
AWS (región Chile o Sudamérica), socio AWS Partner Network (APN), regiones primaria y
secundaria (RT-03.01). \\
\addlinespace
Computación de aplicación &
Microservicios sobre AWS Fargate y Amazon ECS, con despliegue \emph{canary} y
azul-verde (RT-04.06/04.07). \\
\addlinespace
Base transaccional &
Motor relacional con soporte ACID y alta disponibilidad multi-AZ (RT-05.02), separado del
almacén analítico. \\
\addlinespace
Base de datos móvil local &
SQLCipher (AES-256) con almacenamiento encriptado en disco para operación
\emph{offline-first}. \\
\addlinespace
Datos analíticos &
Almacén analítico para OTIF, \emph{fill rate} y costo de servir (RT-05.05/05.29). \\
\addlinespace
Ingesta de telemetría IoT &
Ingesta masiva en tiempo real (Amazon Kinesis) para sensores de temperatura de cadena de
frío. \\
\addlinespace
Seguridad &
TLS 1.3, WAF gestionado, cifrado en reposo y a nivel de campo AES-256, gestor de llaves
(KMS) (RT-11.08/11.10). \\
\addlinespace
Observabilidad &
APM y observabilidad (Datadog/CloudWatch) integrados al NOC 24/7. \\
\bottomrule
\end{tabularx}
\end{table}

\section{Equipamiento a proveer (hardware y software)}

\begin{itemize}
  \item \textbf{Terminales de terreno:} dispositivos móviles \emph{rugged} aptos para
        operatoria a $-$22\,°C, lluvia y una sola mano, con pantalla de alto contraste y
        funcionalidad \emph{offline-first} (RT-17.01; Restricción N° 7).
  \item \textbf{Hojas de \emph{picking} digital:} dispositivos con pantalla operable con
        guantes dentro de la cámara de congelado.
  \item \textbf{Termógrafos vehiculares inalámbricos:} conectados al motor de reglas en
        la nube para alertas de excursión térmica en tiempo real.
  \item \textbf{Sensores IoT de cámara:} captura automática y continua de temperatura.
  \item \textbf{Equipamiento del recinto técnico:} racks independientes de servidores y
        comunicaciones (RT-06.05), fuentes de poder redundantes (RT-08.04), conmutadores
        de núcleo y balanceadores en alta disponibilidad (RT-08.03).
\end{itemize}

\section{Data center primario (site principal on-premise)}

El recinto técnico principal se habilita en las instalaciones del CLIENTE (CD Talca),
con un nivel de disponibilidad de infraestructura del 99,95\,\% (RT-06.01):

\begin{itemize}
  \item \textbf{Energía:} UPS con autonomía mínima de 30 minutos a plena carga
        (RT-06.07), generación autónoma de 24 horas continuas (RT-06.08), instalación
        eléctrica independiente según NCh Elec. 2777 (RT-06.09).
  \item \textbf{Climatización:} precisión redundante N+1 con control de temperatura y
        humedad (RT-06.13), monitoreo en línea de temperatura, humedad y agua
        (RT-06.14).
  \item \textbf{Incendios:} detección temprana por aspiración AnaLASER (RT-06.16),
        extinción automática con agente limpio FM-200 (RT-06.17), extintores portátiles y
        monitoreo integrado (RT-06.18/06.19).
  \item \textbf{Seguridad física:} control de acceso biométrico por biometría facial con
        AFIS como respaldo (RT-06.20), bitácora auditable (RT-06.21), estación de
        enrolamiento y acceso de una persona a la vez con verificación previa
        (RT-06.22/06.23), videovigilancia IP con retención de 30 días (RT-06.24).
  \item \textbf{Rutas de comunicaciones:} doble acometida con ingreso al edificio por
        puntos físicos separados (RT-06.32).
\end{itemize}

\section{Data center secundario y recuperación ante desastres}

Se habilita un sitio secundario en dependencias distintas, con replicación en línea para
producción (RT-07.01):

\begin{itemize}
  \item \textbf{Modalidad:} activo-pasivo declarado y justificado frente al costo, al RTO
        comprometido y a la complejidad operacional (RT-07.01), a distancia suficiente
        para no verse afectado por el mismo evento (RT-07.02).
  \item \textbf{RTO/RPO:} RTO máximo de 4 horas y RPO máximo de 15 minutos para servicios
        críticos (RT-07.04), con replicación continua y alertamiento del retraso
        (RT-07.03).
  \item \textbf{Pruebas:} conmutación real al menos dos veces al año con medición del RTO
        y RPO efectivos (RT-07.07).
  \item \textbf{Respaldos 3-2-1-0:} tres copias, dos medios, una fuera de sitio, una
        inmutable, cero errores de verificación (RT-07.09), cifrados en reposo y tránsito
        con clave gestionada independiente (RT-07.10), copias inmutables protegidas contra
        borrado (RT-07.11), prueba de restauración mensual (RT-07.12).
\end{itemize}

\section{Niveles de servicio de infraestructura}

La infraestructura debe sostener los niveles de disponibilidad del Capítulo 7 y el
compromiso contractual del Artículo 78° (transacción de negocio crítica de extremo a
extremo $\geq$ 99,9\,\%).

\vspace{4pt}
{\small\color{lafroxGrato}\pendiente{} Completar el Formulario T-11 con el detalle
especificado: dimensionamiento por componente, especificaciones de implementos, planos
del recinto y cálculo de capacidad.}
